%Schriftart
\usepackage{fontspec}
%\usepackage[utf8]{inputenc}
%\usepackage[T1]{fontenc}
\usepackage[ngerman]{babel}

%Festlegen des Seitenlayouts
\usepackage[onehalfspacing]{setspace} %Zeilenabstand
%alternativ \linespread{1.5}
%BCOR ist der Binderand, DIV definiert den Satzspiegel (siehe KOMA-script Dokumentation S. 30)
\KOMAoptions{BCOR=0.5cm, DIV=11, parskip=half}
%Absätze werden über parskip geregelt. Die Werte sind parskip=true/false/half

%Abstand vor Kapitelbeginn
\RedeclareSectionCommand[
  beforeskip=0pt,
  afterindent=false
]{chapter}

%Schriftart der Überschriften serif in bold
\addtokomafont{disposition}{\rmfamily}
%Schriftart der Überschriften in serif ohne bold
%\setkomafont{disposition}{\normalfont}

%Für das Einfügen eines Glossars: siehe https://en.wikibooks.org/wiki/LaTeX/Glossary

%stützt die Schnittstelle von float und listings zu tocbasic (genutzt von KOMA-script)
\usepackage{scrhack}

%Zum Auskommentieren von größeren Blöcken
\usepackage{comment}
%Als Platzhalter
\usepackage{lipsum}
%Einfachere Verwendung von Anführungszeichen
\newcommand{\qt}[1]{\glqq#1\grqq{}}
%Angabe von Bildquellen
\newcommand{\quelle}[1]{\\\vspace{.1cm}{Quelle: #1}}

\setcounter{tocdepth}{2}  % show subsections


%Bilder
\usepackage{graphicx}
\usepackage{float}
\graphicspath{{./abbildungen/}}
\usepackage{tabularx}


%Einfügen von Unformatiertem typewriter Text
\usepackage{verbatim}
%über mehrere Zeilen mit \begin{Verbatim}[breaklines=true]\end{Verbatim}
\usepackage{fancyvrb}
\usepackage{makecell}

%Alles was das Formelherz begehrt
\usepackage{amstext} 
\usepackage{amsmath}
\usepackage{amssymb}
\usepackage{booktabs}

%Für die Literaturverzeichnisse
\usepackage{csquotes}
\usepackage[
backend=biber,
style=authoryear,
sorting=nyt,
autocite=footnote,
maxcitenames=2,
dashed=false
]{biblatex}

% Format all names as Family, Given
\DeclareNameAlias{sortname}{family-given}

% Separate authors with semicolons instead of commas
\DeclareDelimFormat{multinamedelim}{\addsemicolon\space}
\DeclareDelimFormat{finalnamedelim}{\addspace\bibstring{and}\space} % uses "und" before the last name

\DefineBibliographyStrings{german}{
	andothers = {et\ al.},
	urlseen   = {Abrufdatum: },
}
\addbibresource{lit.bib}

\setlength{\bibhang}{0pt}
\renewbibmacro*{begentry}{%
	\ifkeyword{bild}{}{\ensuremath{\bullet}\space}%
}

\defbibenvironment{bildliste}
{\list{}
	{\setlength{\leftmargin}{3.5cm}%
		\setlength{\itemindent}{-3.5cm}%
		\setlength{\itemsep}{\bibitemsep}%
		\setlength{\parsep}{\bibparsep}}}
{\endlist}
{\item[]\textbf{\printfield{label}:}\addspace}
%Literaturverzeichnis exakt nach ISO 690-3
%Eine Alternative Erstellung von Literaturverzeichnissen über natbib (letzte Version von 2010)
%WICHTIG: bei einem Wechsel müssen alle Temporären Dateien von LaTex gelöscht werden!
%\usepackage{natbib}
%\usepackage{multibib} %Unterteilung in Internet- und Printmedien
%\newcites{online}{Internetquellenverzeichnis} 
%Statt \cite{} muss für Internetquellen nun \citeonline{} verwendet werden!

\DeclareSortingTemplate{bylabel}{
	\sort{\field{label}}
	\sort{\field{sortname}\field{author}\field{editor}\field{translator}}
	\sort{\field{year}}
	\sort{\field{title}}
}

%Anklickbare Links und \url{}
\usepackage[hidelinks]{hyperref}

\usepackage{xurl}


%Zeilenumbrüche in URLs, standard url Paket wird von hyperref geladen, Befehl \url{} bleibt gleich
\usepackage{microtype}[final]

\usepackage{chngcntr}
\counterwithout{footnote}{chapter}
%Einbinden von Quelltext in Python (https://tex.stackexchange.com/questions/83882/how-to-highlight-python-syntax-in-latex-listings-lstinputlistings-command)
% Default fixed font does not support bold face
\DeclareFixedFont{\ttb}{T1}{txtt}{bx}{n}{12} % for bold
\DeclareFixedFont{\ttm}{T1}{txtt}{m}{n}{12}  % for normal

% Custom colors
\usepackage{color}
\definecolor{deepblue}{rgb}{0,0,0.5}
\definecolor{deepred}{rgb}{0.6,0,0}
\definecolor{deepgreen}{rgb}{0,0.5,0}

\usepackage{listings}

% Python style for highlighting
\newcommand\pythonstyle{\lstset{
language=Python,
basicstyle=\ttm,
otherkeywords={self},             % Add keywords here
keywordstyle=\ttb\color{deepblue},
emph={MyClass,__init__},          % Custom highlighting
emphstyle=\ttb\color{deepred},    % Custom highlighting style
stringstyle=\color{deepgreen},
frame=tb,                         % Any extra options here
showstringspaces=false            % 
}}


% Python environment
\lstnewenvironment{python}[1][]
{
\pythonstyle
\lstset{#1}
}
{}

% Python for external files
\newcommand\pythonexternal[2][]{{
\pythonstyle
\lstinputlisting[#1]{#2}}}

% Python for inline
\newcommand\pythoninline[1]{{\pythonstyle\lstinline!#1!}}
